\subsection{The Big Bang as a Maximal Constraint Regime of the \texorpdfstring{$\chi$}{χ} Substrate}
  \label{subsec:big-bang-maximal-constraint}

  The Big Bang is interpreted as the boundary of applicability of effective
  spacetime descriptions.
  In this regime, the density of structural and topological constraints
  within~$\chi$ exceeds the threshold required for stable geometric
  projection: effective notions of spatial distance, temporal duration, and
  causal ordering cease to be well-defined.
  The apparent singular behavior of standard cosmological models reflects
  the extrapolation of geometric descriptions beyond their domain of validity.

  Cosmological evolution is therefore described as the progressive selection
  of admissible configurations through the effective reprojection sequence,
  starting from this maximal constraint regime.
  The Big Bang marks not the origin of spacetime, but the transition beyond
  which spacetime becomes an appropriate effective framework.
  The arrow of time arises from the irreversibility of the reprojection
  sequence under the admissibility constraints of~$\chi$
  (Section~\ref{subsec:monotonicity-and-arrow-of-time}).

  \paragraph{Structural symmetry breaking at the first transition.}
    The initial state~$U_0$ occupies a unique position in the cascade: it has no
    predecessor~$U_{-1}$.
    The admissible neighbourhood~$\mathcal{A}(U_0)$ is therefore unconditioned by
    any prior projective history, making~$U_0$ maximally symmetric --- all
    successor directions are equivalently accessible.
    The first transition $U_0 \to U_1$ is, however, inevitably asymmetric: once
    a specific~$U_1 \in \mathcal{A}(U_0)$ is reached, the next neighbourhood
    $\mathcal{A}(U_1)$ differs generically from~$\mathcal{A}(U_0)$.
    This is not a phase transition driven by an external potential or a
    temperature parameter.
    It is a structural consequence of the evolution principle itself:
    because every transition shifts the projective position and thus the
    local admissible neighbourhood, perfect symmetry and non-trivial evolution
    are mutually exclusive.
    The apparent symmetry-breaking of the early universe is therefore not
    an event that requires explanation; it is the unavoidable first step
    of any evolution founded on a difference-based admissibility principle.
